\documentclass[notes=show]{beamer}
%%%%%%%%%%%%%%%%%%%%%%%%%%%%%%%%%%%%%%%%%%%%%%%%%%%%%%%%%%%%%%%%%%%%%%%%%%%%%%%%%%%%%%%%%%%%%%%%%%%%%%%%%%%%%%%%%%%%%%%%%%%%%%%%%%%%%%%%%%%%%%%%%%%%%%%%%%%%%%%%%%%%%%%%%%%%%%%%%%%%%%%%%%%%%%%%%%%%%%%%%%%%%%%%%%%%%%%%%%%%%%%%%%%%%%%%%%%%%%%%%%%%%%%%%%%%
\AtBeginSubsection[
  {\frame<beamer>{\frametitle{Outline}   
    \tableofcontents[currentsection,currentsubsection]}}%
]%
{
  \frame<beamer>{ 
    \frametitle{Outline}   
    \tableofcontents[currentsection,currentsubsection]}
}
\usepackage{ulem}
\usepackage{pdflscape}
\setbeamertemplate{caption}{\insertcaption}
\usepackage{color,xcolor,ucs}
\usepackage{beamerthemesplit}
\usepackage{graphicx}
\usepackage{amsmath}
\usepackage{adjustbox}
\usepackage{pdfpages} 
\definecolor{mintgreen}{rgb}{0.6, 1.0, 0.6}
\makeatletter
\@addtoreset{subfigure}{framenumber}% subfigure counter resets every frame
\makeatother
\newcommand*{\Scale}[2][4]{\scalebox{#1}{$#2$}}%
\newcommand*{\Resize}[2]{\resizebox{#1}{!}{$#2$}}%
\usepackage{amsfonts}
\usepackage{booktabs}
\usepackage{color,soul}
\usepackage{pdfpages}
\setboolean{@twoside}{false}
\usepackage{amsmath}
\usepackage{pdflscape} 
\usepackage{eurosym}
\usepackage{amstext} % for \text
\DeclareRobustCommand{\officialeuro}{%
  \ifmmode\expandafter\text\fi
  {\fontencoding{U}\fontfamily{eurosym}\selectfont e}}
\usepackage[caption = false]{subfig}
\usepackage{appendixnumberbeamer} 
\usepackage{graphicx}
\usepackage{mathpazo}
\usepackage{hyperref}
\usepackage{multimedia}
\usepackage{graphicx}
\usepackage{multirow}
\usepackage{subfig}
\usepackage{verbatim}
\usepackage{booktabs, multicol, multirow}
\setcounter{MaxMatrixCols}{10}
\input{Macros.tex}
\AtBeginSection[]
{
  \begin{frame}[noframenumbering]
    \frametitle{Outline for section \thesection}
    \tableofcontents[currentsection]
  \end{frame}
}


\usepackage{xcolor,soul}
\renewcommand<>{\hl}[1]{\only#2{\beameroriginal{\hl}}{#1}}

% https://tex.stackexchange.com/questions/41683/why-is-it-that-coloring-in-soul-in-beamer-is-not-visible
\makeatletter
\newcommand\SoulColor{%
  \let\set@color\beamerorig@set@color
  \let\reset@color\beamerorig@reset@color}
\makeatother
\SoulColor
\AtBeginSection{\frame{\sectionpage}}

\newsavebox\mybox
\newenvironment{aquote}[1]
  {\savebox\mybox{#1}\begin{quote}\openautoquote\hspace*{-.7ex}}
  {\unskip\closeautoquote\vspace*{1mm}\signed{\usebox\mybox}\end{quote}}

\newtheorem{proposition}[theorem]{Proposition}
\newtheorem{Assumption}[theorem]{Assumption}
\newenvironment{stepenumerate}{\begin{enumerate}[<+->]}{\end{enumerate}}
\newenvironment{stepitemize}{\begin{itemize}[<+->]}{\end{itemize} }
\newenvironment{stepenumeratewithalert}{\begin{enumerate}[<+-| alert@+>]}{\end{enumerate}}
\newenvironment{stepitemizewithalert}{\begin{itemize}[<+-| alert@+>]}{\end{itemize} }
\usetheme{Madrid}
\makeatletter
\setbeamertemplate{footline}
{
  \leavevmode%
  \hbox{%
  \begin{beamercolorbox}[wd=.333333\paperwidth,ht=2.25ex,dp=1ex,center]{author in head/foot}%
    \usebeamerfont{author in head/foot}\insertshortauthor~~\beamer@ifempty{\insertshortinstitute}{}{(\insertshortinstitute)}
  \end{beamercolorbox}%
  \begin{beamercolorbox}[wd=.333333\paperwidth,ht=2.25ex,dp=1ex,center]{title in head/foot}%
    \usebeamerfont{title in head/foot}\insertshorttitle
  \end{beamercolorbox}%
  \begin{beamercolorbox}[wd=.333333\paperwidth,ht=2.25ex,dp=1ex,right]{date in head/foot}%
    \usebeamerfont{date in head/foot}\insertshortdate{}\hspace*{2em}
%    \insertframenumber{} / \inserttotalframenumber\hspace*{2ex} % DELETED
  \end{beamercolorbox}}%
  \vskip0pt%
}
\makeatother
\usepackage{pdfpages}
\begin{document}
	\title[]{\large Panel Data
	\vspace{0.5cm}
	}
	\author[Raffaele Saggio - Econ 628]{Raffaele Saggio}
	\institute[]{UBC}
	\date{\today \\
}
	\maketitle

\begin{frame}{What is panel data?}
\begin{itemize}
\item By panel data, we label a data configuration in which we have multiple observations for a given (cross-sectional) unit.
\medskip
\item Typical example: individuals observed over multiple periods.
\medskip
\item Increasing availability of panel data sets $\rightarrow$ ``explosion" of microeconometric studies (both structural and reduced form).
\end{itemize}
\end{frame}
%%
\begin{frame}{Why panel data?}
\begin{enumerate}
\item Selection on observables strategy.
\medskip
\item Permits to Disentangle components of variance.
\medskip
\item Leads to Difference-in-Differences Research Design.  
\end{enumerate}
\end{frame}
\begin{frame}
\frametitle{Main References}
\begin{itemize}
\item Manuel Arellano book + free lecture notes online: \url{https://www.cemfi.es/~arellano/} (mostly for topic 1 and 2). Cameron and Trivedi book also is really good.
\medskip
\item Stuff on the Correlated Random Effect model can be found on Chamberlain's (amazing) handbook chapter published in 1984.
\medskip
\item Mostly Harmeless Econometrics for DiD (next slide-deck).  
\end{itemize}
\end{frame}
%%
\section{The Linear Fixed Effects Model}
%%
\begin{frame}{Farming Panel Data}
\begin{itemize}
\item Suppose you have a cross-sectional (structural) regression of the form
\be
y_{i1} = \beta x_{i1} + \eta_{i} + v_{i1}
\ee
\begin{itemize}
\item $y_{i1}\equiv$ log output.
\item $x_{i1}\equiv$ log labor.
\item $\eta_{i}\equiv$ soil quality.
\item $v_{i1}\equiv$ stochastic input outside farmer's control (e.g. weather).
\end{itemize}
\pause
\item We have that $\E[v_{i1} \vert  x_{i1}, \eta_{i}]=0$. 
\medskip
\item Suppose that $\eta_{i}$ is not observed $\rightarrow$ Short regression won't work (unless...).
 \medskip
 \item No IV's are available.
 \pause
 \item But Suppose I give you another ``round" of data
 \be
 y_{i2} = \beta x_{i2} + \eta_{i} + v_{i2}
 \ee
 with $\E[v_{i2} \vert  x_{i2}, \eta_{i}]=0$
\end{itemize}
\end{frame}
%%
\begin{frame}{Farming Panel Data}
\begin{itemize}
\item One can then take \underline{first differences}
\be
 y_{i2}- y_{i1} = \beta_{1}(x_{i2}-x_{i1}) + v_{i2}-v_{i1}
\ee
\item Is the regression above going to identify $\beta_{1}$?
\medskip
\item Which observations drive the identification here?
\medskip
\pause
\item Panel data allowed us to control for unobserved time invariant heterogeneity. 
\medskip
\item Sometimes (but not always!) panel data can lead to convincing research designs.
\end{itemize}
\end{frame}

\begin{frame}[plain]{A Famous Example: Ashenfelter and Krueger (1994, AER)}
\centering
\includegraphics[scale=0.4]{aux/twins}
\\
Twins Days Festival, a Festival in Twinsburg, Ohio
\end{frame}
\begin{frame}
\frametitle{Ashenfelter and Krueger (1994, AER)}
\begin{itemize}
\item AK attended  the 16th Annual Twins Days Festival in Twinsburg, Ohio, in August of 1992
\medskip
\item This is the largest gathering of twins in the world ($\approx$ 3,000 sets of twins).
\medskip
\item AK collected data on 495 separate individuals over age 18.
\medskip
\item Idea: test for returns to college education by comparing identical twins with different levels of education. 
 \medskip
 \item In this empirical application $\eta_{i}$ genetic and other characteristics that are common across twins. 
 \end{itemize}
\end{frame}
%%%%%%%%
\begin{frame}[plain]{Summary Stats}
\centering
\includegraphics[scale=0.4]{aux/ak1}
\end{frame}
%%%%%%%%
\begin{frame}[plain]{Results}
\centering
\includegraphics[scale=0.4]{aux/ak2}
\be
w_{it} = \alpha_{i} + \beta X_{it} + \varepsilon_{it}
\ee
\end{frame}
%%%%%%%%
\begin{frame}[plain]{The linear FE model aka Heterogeneity as Nuisance}

Textbook panel ``fixed effects'' model:
\[
Y_{it}=\alpha_{i}+X_{it}'\beta+u_{it}\mbox{ \ensuremath{\left(\ensuremath{t=1,...,T}\right)}}
\]
\begin{itemize}
\item Interested in $\beta$, view $\left\{ \alpha_{i}\right\} $
as ``nuisance''. 
\bigskip
\item We are unwilling to impose any restriction on the distribution of $(\alpha_{i} \vert X_{it})$ $\rightarrow$ Can we still identify $\beta$?
\pause
\bigskip
\item ``Fixed Effects" approach: do not impose any assumption on $G(\alpha_{i} \vert X_{i,1}, \hdots  X_{i,T})$.
\bigskip
\item ``Random Effects" approach: impose some restrictions on $G(\alpha_{i} \vert X_{i,1}, \hdots  X_{i,T})$.


\end{itemize}
\end{frame}
\begin{frame}
\frametitle{The Linear FE Model}
\begin{itemize}
\item {\footnotesize{}Three forms of exogeneity (can weaken to non-correlation):}{\footnotesize\par}
\end{itemize}
\begin{enumerate}
\item {\footnotesize{}\emph{{\footnotesize{}Strict}} exogeneity
\[
\mathbb{E}\left[u_{it}|\alpha_{i},X_{i1},...,X_{iT}\right]=0
\]
}{\footnotesize\par}
\item {\footnotesize{}\emph{{\footnotesize{}Sequential}} exogeneity
\[
\mathbb{E}\left[u_{it}|\alpha_{i},X_{i1},...,X_{it}\right]=0
\]
}{\footnotesize\par}
\item {\footnotesize{}\emph{{\footnotesize{}Contemporaneous}} exogeneity
\[
\mathbb{E}\left[u_{it}|\alpha_{i},X_{it}\right]=0
\]
}{\footnotesize\par}
\end{enumerate}
\end{frame}

\begin{frame}[plain]{Within transformation}

\begin{itemize}
\item Subtract $ith$ specific means
\[
Y_{it}-\bar{Y}_{i}=\left(X_{it}-\bar{X}_{i}\right)'\beta+u_{it}-\bar{u}_{i}
\]
\item Does this ring any bell?
\pause
\medskip
\item Estimating the equation above with OLS yields the within-group of fixed effect estimator of $\beta$ (aka the least-square-dummy-variable LSDV estimator).
\medskip
\pause
\item For OLS to be consistent we need
\[
\mathbb{E}\left[\left(u_{it}-\bar{u}_{i}\right)X_{it}\right]=0
\]
\item Strict exogeneity implies these conditions hold
\item But somewhat restrictive b/c $X_{it}$ can't depend on future or past
errors

\begin{itemize}
\item violated in simple forward looking models w/ learning (Mundlak, 1961)
\item see also agriculture example from beginning of class. 
\end{itemize}
\end{itemize}
\end{frame}

\begin{frame}[plain]{First differenced estimation}

\begin{itemize}
\item Alternative strategy is to estimate in 1st differences
\[
Y_{it}-Y_{it-1}=\left(X_{it}-X_{it-1}\right)'\beta+u_{it}-u_{it-1}
\]
\item Here we need only a differenced form of contemporaneous exogeneity
\[
\mathbb{E}\left[\left(u_{it}-u_{it-1}\right)\left(X_{it}-X_{it-1}\right)\right]=0
\]
\end{itemize}

\pause{}
\begin{itemize}
\item However this could still be violated if $X_{it}$ depends on $u_{it-1}$

\begin{itemize}
\item Could always solve by instrumenting $\left(X_{it}-X_{it-1}\right)$
with $X_{it-1}$ (under sequential exog)
\end{itemize}
\end{itemize}

\pause{}
\begin{itemize}
\item What if we run OLS on stacked differences? 
\medskip
\item Can prove that Feasible GLS in stacked differences model is the within group fixed effects estimator, see section 1.2 of \url{https://www.cemfi.es/~arellano/static-panels-class-note.pdf}
\end{itemize}
\end{frame}
%%%%%%%%%%%
\begin{frame}{Correlated Random Effects }

\begin{itemize}
\item Suppose \emph{balanced} panel w/ $T$ time periods

\begin{itemize}
\item See Wooldridge (2009) for extension to unbalanced case
\end{itemize}
\item Linear ``fixed effects'' panel model revisited
\[
Y_{it}=\alpha_{i}+\beta X_{it}+u_{it}
\]
\[
\mathbb{E}^{*}\left[u_{is}|X_{it},\alpha_{i},\right]=0,\mbox{ \ensuremath{\forall s\in\left\{ 1,...,T\right\} }}
\]
where $\mathbb{E}^{*}$ denotes linear projection operator
\item Object of interest is $\beta$, $\left\{ \alpha_{i}\right\} $ are
``nuisance''
\end{itemize}

\pause{}
\begin{itemize}
\item Chamberlain (1982): 

\begin{itemize}
\item $\beta$ is heavily over-identified based upon moments that are estimable
at conventional rates
\item Can use to improve efficiency \& test model
\end{itemize}
\end{itemize}
\end{frame}

\begin{frame}[plain]{The power of thinking ``wide''}

\begin{itemize}
\item Collect observables $\mathbf{X}_{i}=\left[X_{i1},X_{i2},...,X_{iT}\right]'$
and $\mathbf{Y}_{i}=\left[Y_{i1},Y_{i2},...,Y_{iT}\right]'$. 
\pause{}
\begin{itemize}
\item Now Treat $\alpha_{i}$ as ``omitted variable'' and think about its projection on $\mathbf{X}_{i}$ :
\[
\mathbb{E}^{*}\left[\alpha_{i}|\mathbf{X}_{i}\right]=\eta+\lambda_{1}X_{i1}+\lambda_{2}X_{i2}+...+\lambda_{T}X_{iT}\equiv z_{i}
\]
\item Basic idea: need to control for $z_{i}$ to remove OVB!
\end{itemize}
\end{itemize}
\end{frame}

\begin{frame}[plain]{Estimating $z_{i}$}

\begin{itemize}
\item Chamberlain (1982): can do better than Mundlak by imposing a true
restriction on $\mathbb{E}^{*}\left[\alpha_{i}|\mathbf{X}_{i}\right]$ 
\item Get $z_{i}$ from system of reduced form regressions:{\small{}
\begin{eqnarray*}
\mathbb{E}^{*}\left[Y_{it}|\mathbf{X}_{i}\right] & = & \mathbb{E}^{*}\left[\alpha_{i}|\mathbf{X}_{i}\right]+\gamma_{t}+\beta X_{it}\\
 & = & \eta+\gamma_{t}+\lambda_{1}X_{i1}+...+\left(\lambda_{t}+\beta\right)X_{it}+...+\lambda_{T}X_{iT}\\
 & = & \eta+\gamma_{t}+\pi_{t1}X_{i1}+...+\pi_{tt}X_{it}+...+\pi_{tT}X_{iT}
\end{eqnarray*}
}{\small\par}
\end{itemize}

\pause{}
\begin{itemize}
\item Note: we've rewritten a problem with $N$ nuisance parameters $\left\{ \alpha_{i}\right\} _{i=1}^{N}$
as one with $T<<N$ nuisance params $\left\{ \lambda_{t}\right\} _{t=1}^{T}$
\item We can then use CMD to exploit model restrictions
\end{itemize}
\end{frame}

\begin{frame}[plain]{CRE Identification}

\begin{itemize}
\item Collect $\left\{ \pi_{tt'}\right\} _{t=1,t'=1}^{T,T}$ into $T\times T$
matrix $\boldsymbol{\Pi}$ of RF coefficients
\item FE model implies:
\[
\boldsymbol{\Pi}=\beta\boldsymbol{I}_{T}+\boldsymbol{1}_{T}\boldsymbol{\lambda}'
\]
\item Order condition holds when $T\geq2$: 

\begin{itemize}
\item $T^{2}$ RF moments
\item Only $T+1$ structural unknowns
\end{itemize}
\item Rank condition: generically satisfied here
\item CMD estimation $\Rightarrow$ $\chi^{2}$ goodness of fit
\end{itemize}
\end{frame}

\begin{frame}[plain]{Example ($T=2$)}

\begin{eqnarray*}
\boldsymbol{\Pi} & = & \left[\begin{array}{cc}
\pi_{11} & \pi_{12}\\
\pi_{21} & \pi_{22}
\end{array}\right]\\
 & = & \left[\begin{array}{cc}
\beta+\lambda_{1} & \lambda_{2}\\
\lambda_{1} & \beta+\lambda_{2}
\end{array}\right]
\end{eqnarray*}

\begin{itemize}
\item $\beta$ is over-identified:
\[
\beta=\pi_{11}-\pi_{21}=\pi_{22}-\pi_{12}
\]
\end{itemize}

\pause{}
\begin{itemize}
\item Equal and opposite restriction on $\boldsymbol{X}_{i}$:
\[
\Delta Y_{it}=\Delta\gamma_{t}+\beta X_{it}-\beta X_{it-1}+\Delta u_{it}
\]
\item Is violation interesting? 

\begin{itemize}
\item Symmetry: joining a union same as leaving a union?
\end{itemize}
\end{itemize}
\end{frame}

\begin{frame}[plain]{Interpretation issues}

\begin{itemize}
\item Allowing dynamics changes the test. Following model is just-identified
with $T=2$:
\[
Y_{it}=\alpha_{i}+\gamma_{t}+\beta_{1}X_{it}+\beta_{2}X_{it-1}+u_{it}
\]
\item For $T>2$, CMD is testing stationarity of coefficients + restrictions
on dynamics
\end{itemize}
\end{frame}
\begin{frame}{Why panel data?}
\begin{enumerate}
\item Natural selection on observables strategy.
\medskip
\item Disentangle components of variance, study transition and dynamics.  
\end{enumerate}
\pause
\bigskip
We went over point 1. Now we will examine topics related to point 2. Non-Linear FE models are presented in the appendix.  
\end{frame}
%%
\section{Dynamic Panel Data}
\begin{frame}{Dynamics in Panel Data}
\begin{itemize}
\item Often interested in analyzing dynamics, e.g.  separate out permanent from transitory components in consumption decisions. 
\medskip
\item Several examples: Hall and Mishkin (1982), Abowd and Card (1989) and Blundell, Pistaferri, and Preston (2008). 
\medskip
\item Idea: by plotting time series of a given outcome we can infer underlying ``structure" and rule out alternative models.
\end{itemize}
\end{frame}
%%
\begin{frame}{Motivating Example: Genuine Dynamics vs. Unobserved Heterogeneity}
\begin{itemize}
\item First model

\be
y_{it} = \alpha_{i} + v_{it}\qquad t=1,2
\ee
where $\Var(\alpha_{i})=\sigma_{\alpha}^{2}$ and $\Var(v_{it})=\sigma^{2}_{v}$
\end{itemize}
\begin{itemize}
\item Competing model

\be
y_{it} = \alpha + v_{it}\qquad t=1,2 
\ee
\be
v_{it} = \rho v_{it-1} + \varepsilon_{it}\qquad t=1,2 
\ee
where $\vert \rho \vert<1$, $v_{it}\sim_{iid}(0;\sigma^{2}_{v})$; $\varepsilon_{it}\sim_{iid}(0;\sigma^{2}_{\varepsilon})$
\pause
\item Can we distinguish between the two models when $T=2$?
\end{itemize}
\end{frame}
\begin{frame}{Motivating Example: Genuine Dynamics vs. Unobserved Heterogeneity}
\begin{itemize}
\item Can we distinguish between the two models when $T=3$?
\pause
\item Yes! Think about the following moment condition:
\be
\dfrac{\Cov(y_{i1},y_{i2})}{\Cov(y_{i1},y_{i3})} = ?
\ee
\pause
\item How to estimate these models? Usual answer is CMD which provides useful over-id tests as we have seen.
\end{itemize}
\end{frame}
\begin{frame}[plain]{The linear dynamic panel model}

\[
Y_{it}=\alpha_{i}+\delta Y_{it-1}+X_{it}'\beta+\varepsilon_{it}
\]
\[
\mathbb{E}\left[\varepsilon_{it}|\alpha_{i},Y_{it-1},X_{it}\right]=0
\]

\begin{itemize}
\item What is this model saying? 
\end{itemize}
\end{frame}

\begin{frame}
\frametitle{Dynamic Panel Data $\rightarrow$ Restriction on dynamic effects of $X$}
\begin{itemize}
\item By repeated substitution:  
\be
Y_{it} = \dfrac{\alpha_{i}}{1-\delta}+\beta\sum_{k=0}^{\infty}\delta^{k}X_{it-k}+\sum_{k=0}^{\infty}\delta^{k}\varepsilon_{it}
\ee
\item Effect of $X_{it}$ on $Y_{it}$ $\rightarrow$ $\beta$
\item Effect of $X_{it-1}$ on $Y_{it}$ $\rightarrow$ $\beta \delta$.
\item Etc...
\item Dynamic effects dampen geometrically (this is clearly restrictive!).
\end{itemize}
\end{frame}

\begin{frame}[plain]{Estimation}

\begin{itemize}
\item Putting in unit dummies equivalent to demeaning
\[
Y_{it}-\bar{Y}_{i}=\delta\left(Y_{it-1}-\bar{Y}_{i}\right)+\left(X_{it}-\bar{X}_{i}\right)'\beta+\varepsilon_{it}-\bar{\varepsilon}_{i}
\]
\item Nickell (1981) bias: $\mathbb{E}\left[\left(Y_{it-1}-\bar{Y}_{i}\right)\left(\varepsilon_{it}-\bar{\varepsilon}_{i}\right)\right]\neq0$!
\end{itemize}

\pause{}
\begin{itemize}
\item Arellano-Bond (1991): first difference to get
\[
\Delta Y_{it}=\delta\Delta Y_{it-1}+\Delta X_{it}'\beta+\Delta\varepsilon_{it}
\]

\begin{itemize}
\item Instrument $\Delta Y_{it-1}$ using lags of $X_{it}$ and $\left(Y_{i1},...,Y_{it-2}\right)$
\end{itemize}
\end{itemize}

\pause{}
\begin{itemize}
\item Credibility: $\left(Y_{i1},...,Y_{it-2}\right)$ only valid IV if
$\varepsilon_{it}$ serially uncorrelated
\item Statistical problem: 2SLS poorly behaved w/ lots of instruments
\end{itemize}
\end{frame}
%%
\begin{frame}
\centering
\huge{Non-Linear Panel Data Models}
\end{frame}

\begin{frame}[plain]{Panel Probit}

\[
Y_{it}^{*}=\alpha_{i}+\gamma_{t}+\beta X_{it}+\varepsilon_{it}
\]
\[
\varepsilon_{it}\sim N\left(0,1\right)
\]
\[
Y_{it}=1\left\{ Y_{it}^{*}>0\right\} 
\]

\begin{itemize}
\item Heckman (1981): incidental parameters bias small when $T>8$
\item Fernandez-Val (2009) / Fernandez-Val and Weidner (2016): Bias corrections
for FE probit

\begin{itemize}
\item Implemented in Stata w/ ``probitfe'' package
\end{itemize}
\item Alternate approach: Correlated random effects Probit (Chamberlain,
1980)
\end{itemize}
\end{frame}

\begin{frame}[plain]{CRE Probit Model}

\begin{itemize}
\item Restrict conditional \emph{distribution} of unobserved effect
\begin{eqnarray*}
\alpha_{i} & = & \eta+\lambda_{1}X_{i1}+\lambda_{2}X_{i2}+...+\lambda_{T}X_{iT}+v_{i}\\
 & \equiv & z_{i}+v_{i}
\end{eqnarray*}
\[
v_{i}|\boldsymbol{X}_{i},\gamma_{t}\sim N\left(0,\sigma_{v}^{2}\right)
\]
\item Now we have
\begin{eqnarray*}
Y_{it}^{*} & = & \eta+\gamma_{t}+\lambda_{1}X_{i1}+...+\left(\beta+\lambda_{t}\right)X_{it}+...+\lambda_{T}X_{iT}+\varepsilon_{it}\\
 & = & \dot{\gamma}_{t}+\sum_{j=1}^{T}\pi_{tj}X_{ij}+u_{it}
\end{eqnarray*}
where $u_{it}\sim N\left(0,1+\sigma_{v}^{2}\right)$ and $\dot{\gamma}_{t}=\eta+\gamma_{t}$.
\item Can fit this via standard RE Probit routines to get $\left(\sigma_{v}^{2},\Pi\right)$
\end{itemize}
\end{frame}

\begin{frame}[plain]{A limited information approach}

\begin{itemize}
\item Can also just ignore random effect and use the fact that 
\begin{eqnarray*}
\Pr\left(Y_{it}=1|\mathbf{X}_{i}\right) & = & \Phi\left(\frac{\dot{\gamma}_{t}+\sum_{j=1}^{T}\pi_{tj}X_{ij}}{\sqrt{1+\sigma_{v}^{2}}}\right)\\
 & = & \Phi\left(\tilde{\gamma}_{t}+\sum_{j=1}^{T}\tilde{\pi}_{tj}X_{ij}\right)
\end{eqnarray*}
which can be fit via a conventional Probit
\end{itemize}

\pause{}
\begin{itemize}
\item Recover $\tilde{\beta}=\frac{\beta}{\sqrt{1+\sigma_{v}^{2}}}$ in
second step from
\[
\boldsymbol{\tilde{\Pi}}=\tilde{\beta}\boldsymbol{I}_{T}+\boldsymbol{1}_{T}\boldsymbol{\tilde{\lambda}}'
\]
\item Example ($T=2$):
\[
\boldsymbol{\tilde{\Pi}}=\frac{1}{\sqrt{1+\sigma_{v}^{2}}}\left[\begin{array}{cc}
\beta+\lambda_{1} & \lambda_{2}\\
\lambda_{1} & \beta+\lambda_{2}
\end{array}\right]
\]
\end{itemize}
\end{frame}

\begin{frame}[plain]{Which parameters are interesting?}

\begin{itemize}
\item Marginal effect at at given $\alpha_{i}$ depends on $\beta$:
\[
\frac{\partial}{\partial x}\Pr\left(Y_{it}=1|X_{it}=x,\gamma_{t},\alpha_{i}\right)=\beta\phi\left(\beta X_{it}+\gamma_{t}+\alpha_{i}\right)
\]
But which $\alpha_{i}$ to use in evaluating MFX?
\end{itemize}

\pause{}
\begin{itemize}
\item Average partial effect (APE) depends on $\tilde{\beta}$:{\small{}
\begin{eqnarray*}
\mathbb{E}\left[\frac{\partial}{\partial X_{it}}\Pr\left(Y_{it}=1|X_{it},\alpha_{i},\gamma_{t}\right)\right] & = & \mathbb{E}\left[\frac{\partial}{\partial X_{it}}\Pr\left(Y_{it}=1|X_{it},\dot{\gamma}_{t},z_{i},v_{i}\right)\right]\\
 & = & \mathbb{E}\left[\frac{\partial}{\partial X_{it}}\Phi\left(\frac{\beta X_{it}+z_{i}+\dot{\gamma}_{t}}{\sqrt{1+\sigma_{v}^{2}}}\right)\right]\\
 & = & \frac{\beta}{\sqrt{1+\sigma_{v}^{2}}}\mathbb{E}\left[\phi\left(\frac{\beta X_{it}+z_{i}+\dot{\gamma}_{t}}{\sqrt{1+\sigma_{v}^{2}}}\right)\right]\\
 & = & \tilde{\beta}\mathbb{E}\left[\phi\left(\tilde{\beta}X_{it}+\tilde{z}_{i}+\tilde{\gamma}_{t}\right)\right]
\end{eqnarray*}
}{\small\par}
\end{itemize}
\end{frame}

\begin{frame}[plain]{Fixed Effects Logit}

\begin{itemize}
\item Suppose now that error is logit so that
\[
\Pr\left(Y_{it}=1|\boldsymbol{X}_{i},\alpha_{i}\right)=\Lambda\left(X_{it}'\beta+\alpha_{i}\right)
\]
where $\Lambda\left(x\right)=\left[1+\exp\left(-x\right)\right]^{-1}$
\end{itemize}

\pause{}
\begin{itemize}
\item Chamberlain (1980): use sufficient statistic $n_{i}=\sum_{t=1}^{T}Y_{it}$
to remove influence of nuisance effects
\[
\Pr\left(Y_{it}=1|\boldsymbol{X}_{i},\alpha_{i}\right)=\Pr\left(Y_{it}=1|\boldsymbol{X}_{i},n_{i}\right)
\]
\end{itemize}
\end{frame}

\begin{frame}[plain]{Example ($T=2$)}

Logit functional form yields expression analogous to 1st difference:{\scriptsize{}
\begin{eqnarray*}
\Pr\left(Y_{i2}=1|\boldsymbol{X}_{i},\alpha_{i},n_{i}=1\right) & = & \frac{\Pr\left(Y_{i2}=1,n_{i}=1|\boldsymbol{X}_{i},\alpha_{i}\right)}{\Pr\left(n_{i}=1|\boldsymbol{X}_{i},\alpha_{i}\right)}\\
 & = & \frac{\Pr\left(Y_{i2}=1,Y_{i1}=0|\boldsymbol{X}_{i},\alpha_{i}\right)}{\Pr\left(Y_{i2}=1,Y_{i1}=0|\boldsymbol{X}_{i},\alpha_{i}\right)+\Pr\left(Y_{i2}=0,Y_{i1}=1|\boldsymbol{X}_{i},\alpha_{i}\right)}\\
 & = & \frac{\Lambda\left(X_{i2}'\beta+\alpha_{i}\right)\left[1-\Lambda\left(X_{i1}'\beta+\alpha_{i}\right)\right]}{\sum_{\left(t,t'\right)\in\left\{ \left\{ 1,2\right\} ,\left\{ 2,1\right\} \right\} }\Lambda\left(X_{it'}'\beta+\alpha_{i}\right)\left[1-\Lambda\left(X_{it}'\beta+\alpha_{i}\right)\right]}\\
 & = & \Lambda\left(\left(X_{i2}-X_{i1}\right)'\beta\right)
\end{eqnarray*}
}{\scriptsize\par}
\begin{itemize}
\item ML $\Rightarrow$ ``conditional fixed effects'' estimation
\item Interpretation: for \emph{any} $\alpha_{i}$, logit $\beta$ has
interpretation as effect on log odds that $Y_{it}=1$
\end{itemize}
\end{frame}
\begin{frame}[plain]{State Dependence / Heterogeneity}

\begin{itemize}
\item Question: does work experience raise the probability of being employed?

\begin{itemize}
\item heterogeneity: skilled workers more likely to be employed this period
and in the past
\item state dependence: working today \emph{causes} work tomorrow
\end{itemize}

\pause{}
\item Ideal design: randomly assign employment today

\begin{itemize}
\item Or instrument employment today using randomly assigned offer
\end{itemize}
\item Panel literature: use observational variation but control for influence
of time-invariant heterogeneity
\end{itemize}
\end{frame}

\begin{frame}[plain]{State dependent binary choice}

Suppose $Y_{it}$ is indicator for working and $Y_{it}^{*}$ net utility
of working:
\[
Y_{it}^{*}=\alpha_{i}+\gamma Y_{it-1}+\beta X_{it}+\varepsilon_{it}
\]
\[
Y_{it}=1\left\{ Y_{it}^{*}>0\right\} 
\]

\begin{itemize}
\item Note: $Y_{it}^{*}$ depends on $Y_{it-1}^{*}$ only thru $Y_{it-1}$

\begin{itemize}
\item $\gamma$ gives effect of work experience on net utility of working
\end{itemize}
\item Incidental parameters problem when $T$ short. Usual solutions

\begin{itemize}
\item Correlated random effects
\item Conditional fixed effects
\end{itemize}
\end{itemize}
\end{frame}

\begin{frame}[plain]{Random Effects Approach}

\begin{itemize}
\item Likelihood{\footnotesize{}
\begin{eqnarray*}
\Pr\left(Y_{iT},Y_{iT-1},...,Y_{i0}|\boldsymbol{X}_{i}\right) & = & \int\Pr\left(Y_{iT},Y_{iT-1},...,Y_{i0}|\boldsymbol{X}_{i},\alpha\right)dG\left(\alpha|\boldsymbol{X}_{i}\right)\\
 & = & \int\prod_{t=1}^{T}\underbrace{\Pr\left(Y_{it}|Y_{it-1},\boldsymbol{X}_{i},\alpha\right)}_{\mbox{choice model}}\underbrace{\Pr\left(Y_{i0}|\boldsymbol{X}_{i},\alpha\right)}_{\mbox{initial}\mbox{ }\mbox{cond.}}\\
 &  & \times\underbrace{dG\left(\alpha|\boldsymbol{X}_{i}\right)}_{\mbox{heterogeneity}}
\end{eqnarray*}
}{\footnotesize\par}
\end{itemize}

\pause{}
\begin{itemize}
\item Probit choice model: $\Pr\left(Y_{it}=1|Y_{it-1},\boldsymbol{X}_{i},\alpha\right)=\Phi\left(\alpha+\gamma Y_{it-1}+\beta X_{it}\right)$
\item ``Correlated'' normal mixing: $dG\left(\alpha|\boldsymbol{X}_{i}\right)=\frac{1}{\sigma}\phi\left(\frac{\lambda\bar{X}_{i}-\alpha}{\sigma}\right)d\alpha$
\item Initial conditions: 

\begin{itemize}
\item Ignorable when everyone has same $Y_{i0}$ (e.g. sampling kids before
they enter labor market)
\item Ad hoc (but flexible) modeling: $\Pr\left(Y_{i0}|\boldsymbol{X}_{i},\alpha\right)=\Phi\left(\psi\alpha+\phi X_{i0}\right)$
(Heckman, 1981)
\end{itemize}
\end{itemize}
\end{frame}

\begin{frame}[plain]{Initial conditions: alternate factorization}

Wooldridge (2005): condition on 1st period obs{\scriptsize{}
\begin{eqnarray*}
\Pr\left(Y_{iT},Y_{iT-1},...,Y_{i1}|\boldsymbol{X}_{i},Y_{i0}\right) & = & \int\underbrace{\Pr\left(Y_{iT},Y_{iT-1},...,Y_{i1}|\boldsymbol{X}_{i},Y_{i0},\alpha\right)}_{\mbox{choice model}}\\
 &  & \times\underbrace{dG\left(\alpha|\boldsymbol{X}_{i},Y_{i0}\right)}_{\mbox{init cond + het}}
\end{eqnarray*}
}{\scriptsize\par}

\pause{}
\begin{itemize}
\item Normal mixing:
\[
\alpha_{i}=\delta_{0}+\delta_{1}Y_{i0}+\boldsymbol{X}_{i}\zeta+\eta_{i}
\]
\[
\eta_{i}|Y_{i0},\boldsymbol{X}_{i}\sim N\left(0,\sigma_{\eta}\right)
\]
\end{itemize}
\end{frame}

\begin{frame}[plain]{Mass point mixing}

Heckman and Singer (1984): Semi-parametric heterogeneity
\[
dG\left(\alpha|\boldsymbol{X}_{i},Y_{i0}\right)=\sum_{a\in\mathcal{A}_{n}}1\left\{ \alpha=a\right\} \pi_{a}\left(\boldsymbol{X}_{i},Y_{i0}\right)
\]

\begin{itemize}
\item In principle ``sieve space'' $\mathcal{A}_{n}$ grows with $n$
\item Type probabilities
\[
\pi_{a}=\frac{\exp\left(\delta_{0a}+\delta_{1a}Y_{i0}+\boldsymbol{X}_{i}\zeta_{a}\right)}{\sum_{l\in\mathcal{A}_{n}}\exp\left(\delta_{0l}+\delta_{1l}Y_{i0}+\boldsymbol{X}_{i}\zeta_{l}\right)}
\]
\end{itemize}

\pause{}

Bonhomme and Manresa (2015, Ecta): ``Grouped fixed effects''
\begin{itemize}
\item Treat type assignments as parameters to be estimated instead of random
variables
\item Bonhomme, Lamadon, Manresa (2017): requires bias correction when viewed
as sieve
\end{itemize}
\end{frame}

\end{document}
